\section*{6 Conciseness}
In this section we show that $f_0\notin(\End\setminus G)\fcom$. To do so we use a notion called conciseness. Informally, it captures whether a polynomial depends on all variables independent of a change of basis, or a tensor cannot be embedded into a tensor product of smaller spaces.
Given a tensor $f$ in $\C^{m_1}\otimes\C^{m_2}\otimes\cdots\otimes\C^{m_d}$, we associate the following matrices with $f$. For $j\in[d]$, define a matrix $M_f^j$ of dimension $m_j\times(\prod_{i\in[d]\setminus\{j\}}m_i)$ with rows labeled by the standard basis of $\C^{m_j}$, and columns by elements in the Cartesian product
\begin{multline*}
\{\text{standard basis of }\C^{m_1}\}\times\cdots\times\{\text{standard basis of }\C^{m_{j-1}}\}\\
\times\{\text{standard basis of }\C^{m_{j+1}}\}\times\cdots\times\{\text{standard basis of }\C^{m_d}\}.
\end{multline*}
We write the tensor $f$ in the standard basis
\[f=\sum_{\substack{1\le i_1\le m_1\\1\le i_2\le m_2\\\cdots\\1\le i_d\le m_d}}\alpha_{i_1,\ldots,i_d}e_{i_1}\otimes\cdots\otimes e_{i_d}\]
and associate to it the matrix $M_f^j$ whose entry at position $((i_j),(i_1,i_2,\ldots,i_{j-1},i_{j+1},\ldots,i_d))$ is $\alpha_{i_1,\ldots,i_d}$.
\begin{definition}
We say that a tensor $f$ in $\C^{m_1}\otimes\C^{m_2}\otimes\cdots\otimes\C^{m_d}$ is concise if and only if for all $j\in[d]$, $M_f^j$ has full rank.\footnote{When $f$ is viewed as a set-multilinear polynomial (see [36, Section 1.4]), this condition translates to the linear independence of the partial derivatives of $f$. In particular, $M_f^j$ is testing if the partial derivatives of $f$ with respect to the $j$-th block of variables are all linearly independent. This partial derivatives based criterion for testing if a polynomial depends on all the variables, independent of a change of basis, is pretty standard: see, for instance, [23, Corollary5.1.4].}
\end{definition}
As a warm-up exercise we now show that $\fcom$ is concise.
\begin{proposition}$\fcom$ is concise.\end{proposition}
\begin{proof}
We know that $\fcom\in W$. Let us consider the matrix $M^j_{\fcom}$ for some $j\in[d]$. To establish that $M^j_{\fcom}$ has full rank, it suffices to show that rows are linearly independent. In order to show that, we argue that every row is non-zero and every column has at most one non-zero entry. In other words, rows are supported on disjoint sets of columns.
A row of $M^j_{\fcom}$ is labeled by an edge in the $j$-th layer of the ABP $\Gcom$. Recall that only paths that start at a vertex in $V^1$ and end at the corresponding vertex in $V^{d+1}$ contribute to the computation in $\Gcom$. We call such paths valid paths. An entry in $M^j_{\fcom}$ is non-zero iff the corresponding row and column labels form a valid path in $\Gcom$. Thus, it is easily seen that a row is non-zero iff there is a valid path in $\Gcom$ that passes through the edge given by the row label. By the structure of $\Gcom$, in particular that every layer is a complete bipartite graph, we observe that passing through every edge there is some valid path. Hence, we obtain that every row is non-zero.
The second claim now follows from the observation that fixing $d-1$ edges either defines a unique $d$th edge so that these $d$ edges form a valid path, or for these $d-1$ edges there is no such $d$th edge.
\end{proof}
As mentioned in Section 5, to establish $f_0\notin(\End\setminus G)\fcom$ we will show that $f_0$ is concise while any element in $(\End\setminus G)\fcom$ is not.
\begin{lemma}$f_0$ is concise.\end{lemma}
\begin{proof}
Analogous to the proof of Proposition 6, we again show that every row of $M^j_{f_0}$ is non-zero and every column of it has at most one non-zero entry. That is, rows of $M^j_{f_0}$ are supported on disjoint sets of columns.
From the construction of $f_0$ it is seen that a path in $\Gcom$ contributes to the computation of $f_0$ iff it is a valid path that comprises of exactly one parity preserving edge. The second claim of every column having at most one non-zero entry now follows for the same reason as in the proof of Proposition 6.
Before proving the first claim, we recall two assumptions in the construction of $f_0$. The first is that the format $w=(w_1,w_2,\ldots,w_d)$ is such that $w_i\ge2$ for all $i\in[d]$ and the second is that $d$ is odd. To argue that a row is non-zero it suffices to show that a valid path comprising of only one parity preserving edge passes through the edge given by the row level.
Let us consider an arbitrary edge $e$ in $\Gcom$. We have two cases to consider depending on whether it is parity preserving or changing.
\paragraph{Case1.} Suppose $e$ is parity preserving and it belongs to a layer $j\in[d]$. The number of layers on the left of $e$ is $j-1$ and on the right is $d-j$. Since $d$ is odd, these numbers are either both even or both odd. We now argue for the case when they are even (the odd case is analogous). Choose a vertex $v$ in $V^1$ that has the same parity (different in the odd case) as one of the end points of $e$. (Such a choice exists because $w_1\ge2$.) We now claim that there exists a valid path starting at $v$ that passes through $e$ and contains exactly one parity preserving edge. Since $e$ is parity preserving, all edges in the claimed path must be parity changing. We observe that $e$ can be easily extended in both directions using parity changing edges such that the path ends at $\operatorname{corr}(v)$. The existence of parity changing edges at each layer uses the assumption that $w_i\ge2$.
\paragraph{Case2.} Otherwise $e=(a,b)$ is parity changing. Again as before there are two cases based on whether both $j-1$ and $d-j$ are even or odd. Consider the case when they are even (the odd case being analogous). We first assume that $j\ne d$. Choose a vertex $v$ in $V^1$ that has the same parity as $a$. We now construct a valid path from $v$ to $\operatorname{corr}(v)$ that passes through $e$ and contains exactly one parity preserving edge. It is easily seen that there exists a path from $v$ to $a$ using only parity changing edges. We choose a parity preserving outgoing edge incident to $b$. We call its endpoint $v_1$. Since $v_1$ and $v$ have different parities, we can connect $v_1$ to $\operatorname{corr}(v)$ in $V^{d+1}$ using only parity changing edges. Thus we obtain the following valid path $v\to\cdots\to a\to b\to v_1\to\cdots\to\operatorname{corr}(v)$ passing through exactly one parity preserving edge $(b,v_1)$. In the case that $j=d$, choose an incoming parity preserving edge incident on $a$ instead of an outgoing edge on $b$.
\end{proof}
\begin{remark}
We note that if the format $w=(w_1,\ldots,w_d)$ defining $f_0$ is such that for some $j\in[d]$, $w_j=1$, then $f_0$ is not concise. This can be seen as follows.
Let $w_j=1$, and let $v$ denote the unique vertex in $V^j$. Let $e$ be the edge $e=(1,1,j)$. If $j<d$, let $e'$ be the edge $e'=(1,1,j+1)$, otherwise let $e'$ be the edge $e'=(1,1,j-1)$. Both $e,e'$ are parity preserving edges. By construction, every valid path using $e'$ must also use $e$. Hence the corresponding row in the matrix $M^{j+1}_{f_0}$ if $j<d$, and in $M^{j-1}_{f_0}$ otherwise, is zero. Therefore $f_0$ is not concise.
This is an interesting observation, because this is the point where our proof fails for single-(source,sink) ABPs, and this is expected, because Nisan [35] had shown that the set of polynomials computed by such ABPs of format $w$ is a closed set.
\end{remark}
\begin{lemma}Let $f\in(\End\setminus G)\fcom$. Then $f$ is not concise.\end{lemma}
\begin{proof}
This statement is true in very high generality. In our specific case a proof goes as follows. If $f\in(\End\setminus G)\fcom$, then $f=g\fcom$ for some $g\in\End\setminus G$. Let $g=(g_1,\ldots,g_d)$, where $g_i\in\C^{w_iw_{i+1}\times w_iw_{i+1}}$. Since $g\notin G$, at least one of the $g_i$ must be singular. The crucial property is $M^i_{g\fcom}=g_iM^i_{\fcom}$, which finishes the proof.
\end{proof}
